\documentclass[11pt,a4paper]{article}

% --- Codificación e idioma ---
\usepackage[utf8]{inputenc}
\usepackage[T1]{fontenc}
\usepackage[spanish,es-tabla,es-noshorthands]{babel}
\usepackage{lmodern}

% --- Geometría y formato general ---
\usepackage[margin=2.5cm]{geometry}
\usepackage{microtype}
\usepackage{parskip}
\usepackage{enumitem}
\setlist{topsep=2pt,itemsep=2pt,parsep=0pt}

% --- Tablas y diagramas ---
\usepackage{array}
\usepackage{booktabs}
\usepackage{amsmath,amssymb}
\usepackage{tikz}
\usetikzlibrary{positioning,arrows.meta,calc,shapes.geometric}

% --- Cabeceras y numeración ---
\usepackage{fancyhdr}
\pagestyle{fancy}
\fancyhf{}
\fancyhead[L]{\small Bases de Datos II}
\fancyhead[R]{\small Resumen --- XML y JSON}
\fancyfoot[C]{\thepage}
\renewcommand{\headrulewidth}{0.4pt}

% --- Color y código ---
\usepackage{xcolor}
\definecolor{codebg}{RGB}{248,248,248}
\definecolor{codeframe}{RGB}{200,200,200}
\definecolor{codekw}{RGB}{0,0,180}
\definecolor{codestr}{RGB}{163,21,21}
\definecolor{codecmt}{RGB}{0,128,0}
\definecolor{notebg}{RGB}{255,250,230}
\definecolor{noteframe}{RGB}{220,180,80}

\usepackage{listings}
\lstdefinestyle{plain}{
    basicstyle=\ttfamily\footnotesize,
    backgroundcolor=\color{codebg},
    frame=single,
    rulecolor=\color{codeframe},
    breaklines=true,
    columns=fullflexible,
    keepspaces=true,
    xleftmargin=0.4em,
    xrightmargin=0.4em,
    aboveskip=8pt,
    belowskip=8pt
}
\lstset{style=plain}

% --- Cajas de nota ---
\usepackage[most]{tcolorbox}
\newtcolorbox{nota}{colback=notebg,colframe=noteframe,arc=2pt,
    left=6pt,right=6pt,top=4pt,bottom=4pt,boxrule=0.6pt}

% --- Hipervínculos ---
\usepackage[hidelinks,bookmarks=true]{hyperref}

\title{Resumen Teórico --- XML y JSON\\[4pt]
       \large Bases de Datos II --- Práctico N.\textsuperscript{o} 7}
\author{Tomás Rodeghiero}
\date{2026}

\begin{document}

\begin{titlepage}
    \centering
    \vspace*{1.5cm}

    {\large \textbf{Universidad Nacional de Río Cuarto}}\\[6pt]
    {\normalsize Facultad de Ciencias Exactas, Físico-Químicas y Naturales}\\[4pt]
    {\normalsize Departamento de Computación}\\[4pt]
    {\normalsize Licenciatura en Ciencias de la Computación}\\[3cm]

    {\Large \textbf{Bases de Datos II}}\\[1.2cm]

    {\Large \textbf{Resumen Teórico}}\\[6pt]
    {\large XML y JSON --- Bases documentales y MongoDB\\[2pt]
    Práctico N.\textsuperscript{o} 7}\\[3cm]

    \begin{tabular}{rl}
        \textbf{Alumno:} & Tomás Rodeghiero \\[4pt]
        \textbf{Año:}    & 2026 \\
    \end{tabular}

    \vfill
    {\small Síntesis del Teórico 9 (XML y JSON), complementada con la
    parte práctica de MongoDB del Práctico 7.}
\end{titlepage}

\tableofcontents
\newpage

% =====================================================
\section{Datos semiestructurados: idea general}
% =====================================================

Las bases de datos clásicas exigen que cada fila tenga \emph{exactamente}
los mismos campos que el esquema de su tabla. Esto es ideal cuando los
datos están bien definidos, pero genera problemas cuando hay que
intercambiar datos entre aplicaciones diferentes ---cada una con su
propia idea de qué guardar--- o cuando los datos son intrínsecamente
\emph{irregulares} (no todos los productos tienen un ISBN, no todos los
clientes tienen un segundo teléfono, etc.).

Los formatos \textbf{semiestructurados} resuelven este desbalance. Un
documento semiestructurado lleva ``etiquetas'' o ``claves'' que
identifican cada pedazo del dato, y la estructura puede variar de un
documento a otro. Los dos formatos que dominan hoy son:

\begin{itemize}
    \item \textbf{XML}: surgió como lenguaje estándar para intercambio
    de datos entre sistemas. Es verboso pero rico en herramientas
    (esquemas, transformaciones, consultas).
    \item \textbf{JSON}: nació como serialización de objetos JavaScript
    y se volvió el formato dominante por su simpleza. Es la base de las
    APIs REST y de las bases de datos NoSQL documentales (MongoDB).
\end{itemize}

% =====================================================
\section{XML}
% =====================================================

\textbf{XML} (\emph{eXtensible Markup Language}) fue definido por la
World Wide Web Consortium (W3C) como lenguaje genérico para representar
datos estructurados con marcas.

\paragraph{Idea.} Un documento XML está hecho de elementos delimitados
por marcas de apertura y cierre, anidables entre sí:

\begin{lstlisting}
<supermercado>
    <articulo>
        <nro_articulo>10</nro_articulo>
        <descripcion>Harina</descripcion>
        <precio>20.8</precio>
    </articulo>
    <proveedor>
        <nro_proveedor>202</nro_proveedor>
        <nombre>Molinos Canuelas</nombre>
    </proveedor>
</supermercado>
\end{lstlisting}

Las marcas \emph{describen} el contenido: el dato es ``relativamente
autodocumentado''.

\paragraph{Motivación.}
\begin{itemize}
    \item Comunicación de datos entre aplicaciones (bancos,
    procesamiento de órdenes entre empresas).
    \item Base del protocolo SOAP en Web Services.
    \item Texto plano con headers que indican el significado: fácil de
    debuggear y de pasar entre lenguajes.
\end{itemize}

\paragraph{Comparación con datos relacionales.}
\begin{itemize}
    \item \emph{Más ineficiente en almacenamiento}: las marcas se repiten
    en cada elemento, mientras que en una tabla la estructura aparece
    una sola vez en el encabezado.
    \item \emph{Mejor para intercambio}: cada documento es
    autodescriptivo y la estructura puede ser flexible.
\end{itemize}

\subsection{Esquemas: DTD y XML Schema}

Un documento XML \emph{no} requiere asociarse a un esquema, pero los
esquemas son críticos para el intercambio: permiten validar que el
documento recibido es el esperado. Existen dos mecanismos:

\paragraph{DTD (\emph{Document Type Definition}).} La forma original,
hoy obsoleta. Define qué elementos puede haber, qué atributos, qué
subelementos y con qué cardinalidad. No restringe tipos de dato: todo
es \emph{string}.

\begin{lstlisting}
<!ELEMENT element (subelements-specification) >
<!ATTLIST element (attributes) >
\end{lstlisting}

\paragraph{XML Schema.} El estándar actual. Soporta:
\begin{itemize}
    \item tipos de dato (int, string, date, decimal, etc.);
    \item restricciones de mínimos y máximos;
    \item tipos definidos por el usuario;
    \item claves (\texttt{xs:key}) y claves foráneas (\texttt{xs:keyref})
    con \texttt{xs:selector} y \texttt{xs:field};
    \item espacios de nombres (\emph{namespaces}).
\end{itemize}

Está escrito en sintaxis XML (a diferencia de DTD), por lo que se valida
con las mismas herramientas que cualquier otro XML.

Ejemplo de restricción de clave en XML Schema:

\begin{lstlisting}
<xs:key name="claveCuenta">
    <xs:selector xpath="/banco/cuenta"/>
    <xs:field    xpath="nro_cuenta"/>
</xs:key>

<xs:keyref name="fkCuenta" refer="claveCuenta">
    <xs:selector xpath="/banco/titular"/>
    <xs:field    xpath="numero_cuenta"/>
</xs:keyref>
\end{lstlisting}

\subsection{Consultas y transformaciones}

Hay tres lenguajes estándares para trabajar con XML:

\begin{itemize}
    \item \textbf{XPath}: lenguaje de expresiones tipo ``path'' para
    navegar el árbol XML. Es la base sobre la que se construyen los
    otros dos.
    \item \textbf{XQuery}: lenguaje de consulta general para XML
    (análogo a SQL para relacional).
    \item \textbf{XSLT}: lenguaje para transformar XML en otro XML, o
    en HTML.
\end{itemize}

\subsubsection*{XPath: navegación}

Expresiones tipo \texttt{/banco/cuenta[saldo>400]/nro\_cuenta}: bajan
por el árbol seleccionando nodos que cumplen el predicado.

\subsubsection*{XQuery: FLWOR}

XQuery usa la estructura \textbf{FLWOR}: \emph{For, Let, Where, Order by,
Return}. La equivalencia con SQL es:

\begin{center}
\begin{tabular}{ll}
\toprule
XQuery & SQL \\
\midrule
\texttt{for}      & \texttt{from}     \\
\texttt{where}    & \texttt{where}    \\
\texttt{order by} & \texttt{order by} \\
\texttt{return}   & \texttt{select}   \\
\texttt{let}      & (sin equivalencia: variables temporales) \\
\bottomrule
\end{tabular}
\end{center}

Ejemplo: obtener los números de cuenta con saldo $>$ 400:

\begin{lstlisting}
for   $x in /banco-2/cuenta
let   $nroCuenta := $x/@numero-cuenta
where $x/saldo > 400
return <numero-cuenta>{ $nroCuenta }</numero-cuenta>
\end{lstlisting}

\subsubsection*{XSLT: transformaciones}

XSLT (\emph{XML Stylesheet Language Transformations}) usa
\emph{templates} para transformar nodos XML. Cada template hace
\emph{match} sobre un patrón XPath y produce salida con
\texttt{xsl:value-of}:

\begin{lstlisting}
<xsl:template match="/banco-2/cliente">
    <xsl:value-of select="nombre_cliente"/>
</xsl:template>
\end{lstlisting}

Su uso típico es generar HTML a partir de XML, pero sirve para cualquier
conversión XML $\rightarrow$ XML o XML $\rightarrow$ texto.

% =====================================================
\section{JSON}
% =====================================================

\textbf{JSON} (\emph{JavaScript Object Notation}) es un formato liviano
de intercambio de datos. Su éxito está en tres atributos: es muy simple
de interpretar, es independiente del lenguaje y usa convenciones
familiares para programadores de C, Java, JavaScript, Python, etc.

\paragraph{Por qué desplaza a XML.}
\begin{itemize}
    \item Menos verboso (no hay marcas de cierre repetidas).
    \item Mapeo directo a estructuras de datos nativas (objetos,
    arreglos).
    \item Es el formato natural de APIs REST.
    \item Es la base de almacenamiento de las bases NoSQL documentales,
    como MongoDB.
\end{itemize}

\subsection{Estructura de un documento JSON}

JSON está construido sobre dos estructuras:

\begin{itemize}
    \item \textbf{Objeto}: conjunto desordenado de pares
    \texttt{nombre/valor}. Sintaxis: \texttt{\{ "clave": valor, ... \}}.
    \item \textbf{Arreglo}: lista ordenada de valores. Sintaxis:
    \texttt{[ v1, v2, ... ]}.
\end{itemize}

\paragraph{Tipos de valor.} Un valor puede ser:
\begin{itemize}
    \item un \emph{string} entre comillas dobles;
    \item un \emph{número} (entero o decimal);
    \item \texttt{true}, \texttt{false} o \texttt{null};
    \item otro objeto;
    \item un arreglo.
\end{itemize}

Las estructuras se pueden anidar libremente.

\paragraph{Ejemplo.}

\begin{lstlisting}
{
  "nombre":   "Carlos",
  "apellido": "Perez",
  "edad":     45,
  "direccion": {
    "calle":       "San Martin",
    "ciudad":      "Rio Cuarto",
    "codigoPostal": 5800
  },
  "telefonos": [
    { "tipo": "fijo",  "numero": "358-45458412" },
    { "tipo": "movil", "numero": "351-468213744" }
  ]
}
\end{lstlisting}

\subsection{JSON como modelo de datos}

Lo interesante para una base de datos es que JSON soporta directamente:

\begin{itemize}
    \item \textbf{Composición} (objetos embebidos, como
    \texttt{direccion} dentro de la persona): equivale al concepto E/R
    de \emph{atributo compuesto}.
    \item \textbf{Multivalor} (arreglos como \texttt{telefonos}):
    equivale al concepto E/R de \emph{atributo multivaluado}.
    \item \textbf{Heterogeneidad}: dos documentos de la misma colección
    pueden tener distintos campos.
\end{itemize}

Es exactamente lo contrario a la primera forma normal del modelo
relacional. En vez de \emph{romper} estos atributos en tablas separadas,
los modelos documentales los mantienen \emph{juntos} con el resto del
dato.

% =====================================================
\section{Consultas sobre JSON}
% =====================================================

A diferencia de XML, no hay todavía un estándar único de consulta sobre
JSON. La especificación XQuery 3.1 de W3C extiende XQuery a JSON, pero
en la práctica conviven varios lenguajes específicos:

\begin{itemize}
    \item \textbf{JsonPath}: el más usado, inspirado en XPath. Está
    implementado para muchos lenguajes (la cátedra usa Jayway JsonPath
    para Java).
    \item \textbf{JSONiq}: basado en XQuery.
    \item \textbf{JSON Query}: otro lenguaje tipo XPath.
    \item Cada lenguaje de programación tiene además su propio DSL
    (LINQ to JSON en .NET, \texttt{json\_path} en Postgres, etc.).
\end{itemize}

\subsection{JsonPath en detalle}

JsonPath se basa en una colección de \textbf{operadores} y
\textbf{funciones}:

\begin{center}
\begin{tabular}{ll}
\toprule
Operador & Significado \\
\midrule
\texttt{\$}                & Raíz del documento. \\
\texttt{@}                 & Nodo actual (dentro de un predicado). \\
\texttt{*}                 & Comodín. \\
\texttt{..}                & Recorrido en profundidad. \\
\texttt{.<nombre>}         & Acceso a hijo. \\
\texttt{['<n1>','<n2>']}   & Hijos anotados entre corchetes. \\
\texttt{[<i>,<j>]}         & Uno o más índices de arreglo. \\
\texttt{[start:end]}       & Rango de índices. \\
\texttt{[?(<expr>)]}       & Filtro: la expresión devuelve booleano. \\
\bottomrule
\end{tabular}
\end{center}

\paragraph{Funciones de agregación.} \texttt{min()}, \texttt{max()},
\texttt{avg()}, \texttt{length()}.

\paragraph{Operadores de filtro.} \texttt{==}, \texttt{!=}, \texttt{<},
\texttt{<=}, \texttt{>}, \texttt{>=}, \texttt{=$\sim$} (regex),
\texttt{in}, \texttt{nin}, \texttt{subsetof}, \texttt{anyof},
\texttt{noneof}, \texttt{size}, \texttt{empty}.

\subsection{Ejemplos sobre la tienda de libros}

Sobre el JSON clásico de \texttt{store.book[]} con \texttt{bicycle} y
\texttt{expensive} en la raíz:

\begin{center}
\begin{tabular}{ll}
\toprule
JsonPath & Resultado \\
\midrule
\texttt{\$.store.book[*].author}                   & Autores de todos los libros. \\
\texttt{\$..author}                                & Todos los \texttt{author} del documento. \\
\texttt{\$.store.*}                                & Todo lo que la tienda almacena. \\
\texttt{\$.store..price}                           & Precios de todo lo almacenado. \\
\texttt{\$..book[2]}                               & El tercer libro. \\
\texttt{\$..book[-2]}                              & El penúltimo libro. \\
\texttt{\$..book[0,1]}                             & Los primeros dos libros. \\
\texttt{\$..book[?(@.isbn)]}                       & Libros que tienen ISBN. \\
\texttt{\$.store.book[?(@.price < 10)]}            & Libros con precio menor a 10. \\
\texttt{\$..book[?(@.price <= \$['expensive'])]}   & Libros considerados baratos. \\
\texttt{\$..book[?(@.author =$\sim$ /.*REES/i)]}   & Equivalente al \texttt{LIKE} de SQL. \\
\texttt{\$..book.length()}                         & Cantidad de libros. \\
\bottomrule
\end{tabular}
\end{center}

\subsection{Uso desde Java (Jayway)}

\begin{lstlisting}
import com.jayway.jsonpath.JsonPath;
import java.io.File;

File   jsonFile = new File("/tmp/json.txt");
String jsonExp  = "$.store.book[?(@.price > 10)]";

System.out.println(JsonPath.read(jsonFile, jsonExp));
\end{lstlisting}

% =====================================================
\section{Bases de datos documentales: MongoDB}
% =====================================================

Las bases de datos NoSQL orientadas a documentos almacenan JSON (o una
variante binaria como BSON, en el caso de MongoDB) como unidad de dato.
\textbf{MongoDB} es el motor documental más popular y el que usa el
Práctico 7.

\subsection{Conceptos básicos}

\begin{itemize}
    \item \textbf{Base de datos}: contenedor lógico de colecciones.
    \item \textbf{Colección}: grupo de documentos relacionados
    (equivalente a una tabla, pero sin esquema rígido).
    \item \textbf{Documento}: un objeto JSON/BSON con un \texttt{\_id}
    único. Es la unidad mínima de almacenamiento.
    \item \textbf{Esquema}: opcional. Se puede validar con
    \texttt{\$jsonSchema} al crear la colección.
\end{itemize}

\subsection{Operaciones CRUD básicas}

\begin{lstlisting}
// Insert
db.usuarios.insertOne({ nick: "moises", edad: 45, ... });
db.usuarios.insertMany([{...},{...}]);

// Read (find)
db.usuarios.find({ edad: { $gt: 40 } });
db.usuarios.find(
  { edad: { $gt: 40 } },             // filtro
  { _id: 0, nick: 1, edad: 1 }       // proyeccion
);

// Update
db.usuarios.updateOne(
  { nick: "moises" },
  { $set: { edad: 46 } }
);

// Upsert (update o insert si no existe)
db.usuarios.updateOne(
  { nick: "elias" },
  { $set: { ... } },
  { upsert: true }
);

// Delete
db.usuarios.deleteOne({ nick: "moises" });
\end{lstlisting}

\subsection{Operadores de consulta}

Los \emph{filtros} se construyen con operadores tipo \texttt{\$op}:

\begin{itemize}
    \item Comparación: \texttt{\$eq}, \texttt{\$ne}, \texttt{\$gt},
    \texttt{\$gte}, \texttt{\$lt}, \texttt{\$lte}, \texttt{\$in},
    \texttt{\$nin}.
    \item Lógicos: \texttt{\$and}, \texttt{\$or}, \texttt{\$not},
    \texttt{\$nor}.
    \item Existencia: \texttt{\$exists}.
    \item Sobre arreglos: \texttt{\$all}, \texttt{\$size},
    \texttt{\$elemMatch}.
    \item Texto: expresiones regulares en literales o con \texttt{\$regex}.
\end{itemize}

\paragraph{Proyección.} El segundo parámetro de \texttt{find} dice qué
campos incluir (\texttt{1}) o excluir (\texttt{0}). El \texttt{\_id}
viene por defecto incluido y suele excluirse explícitamente.

\subsection{Aggregation framework}

Cuando \texttt{find} no alcanza, MongoDB ofrece un \emph{pipeline} de
etapas. Cada etapa toma los documentos de la anterior y los transforma:

\begin{center}
\begin{tabular}{ll}
\toprule
Etapa & Función \\
\midrule
\texttt{\$match}     & Filtra documentos (como \texttt{find}). \\
\texttt{\$project}   & Selecciona y calcula campos. \\
\texttt{\$group}     & Agrupa por una clave (similar a \texttt{GROUP BY}). \\
\texttt{\$sort}      & Ordena. \\
\texttt{\$limit} / \texttt{\$skip} & Paginación. \\
\texttt{\$unwind}    & Desempaqueta un arreglo en un documento por elemento. \\
\texttt{\$lookup}    & Join lógico con otra colección. \\
\texttt{\$addFields} & Agrega o transforma campos. \\
\bottomrule
\end{tabular}
\end{center}

\paragraph{Ejemplo: \texttt{\$lookup} para hacer un join lógico.}

\begin{lstlisting}
db.noticias.aggregate([
  {
    $lookup: {
      from:         "usuarios",
      localField:   "autor_nick",
      foreignField: "nick",
      as:           "autor"
    }
  },
  { $unwind: "$autor" },
  { $match: { "autor.nyape": /gonzalez$/i } },
  { $sort:  { fecha: 1 } }
]);
\end{lstlisting}

\subsection{Decisiones de modelado documental}

\paragraph{Embedding vs Referencing.} La decisión clave al modelar en
MongoDB:

\begin{itemize}
    \item \textbf{Embedding}: anidar un documento dentro de otro. Ideal
    cuando el documento embebido es \emph{parte} del padre y la
    consulta dominante los pide juntos. Ejemplo: dirección dentro del
    usuario.
    \item \textbf{Referencing}: guardar una referencia (clave) a otra
    colección. Ideal cuando las entidades tienen ciclos de vida
    independientes o cuando el embebido crecería sin tope. Ejemplo:
    comentarios de una noticia (pueden ser muchísimos y se agregan a lo
    largo del tiempo).
\end{itemize}

\paragraph{Sin foreign keys.} MongoDB \emph{no} valida referencias entre
colecciones. La consistencia referencial queda del lado de la
aplicación o de los validadores manuales.

\paragraph{Índices.} Crear índices únicos sobre las claves de negocio
(\texttt{nick}, \texttt{codigo}) e índices simples sobre los campos por
los que se filtra o se hace \texttt{\$lookup} es indispensable para que
las consultas sean eficientes a escala.

\subsection{Validación con \texttt{\$jsonSchema}}

Aunque MongoDB no exige esquema, se le puede dar uno al crear la
colección:

\begin{lstlisting}
db.createCollection("usuarios", {
  validator: { $jsonSchema: {
    bsonType: "object",
    required: ["nick","nyape","edad","direccion","telefonos"],
    properties: {
      nick:      { bsonType: "string", minLength: 1 },
      edad:      { bsonType: "int",    minimum:   0 },
      direccion: { bsonType: "object",
                   required: ["calle","altura"],
                   properties: {
                     calle:  { bsonType: "string", minLength: 1 },
                     altura: { bsonType: "int",    minimum:   0 }
                   }},
      telefonos: { bsonType: "array",  minItems: 1,
                   items: { bsonType: "string", minLength: 1 }}
    }
  }}
});
\end{lstlisting}

Da garantías de tipos sin sacrificar la flexibilidad documental.

% =====================================================
\section{XML vs JSON: comparación}
% =====================================================

\begin{center}
\begin{tabular}{lll}
\toprule
Dimensión & XML & JSON \\
\midrule
Verbosidad         & Alta (etiquetas duplicadas)     & Baja \\
Tipos de dato      & Todo string en DTD; tipados en XML Schema & Nativos: number, bool, null \\
Esquemas           & DTD / XML Schema (maduros)      & JSON Schema (más reciente) \\
Consultas          & XPath, XQuery (estándares)      & JsonPath, JSONiq (no estandarizados) \\
Transformaciones   & XSLT                             & Ad-hoc (mapeo en código) \\
Uso típico         & Documentos textuales, SOAP       & APIs REST, NoSQL, frontend \\
Bases que lo usan  & XML nativos, Oracle XML DB       & MongoDB, CouchDB, Postgres JSONB \\
\bottomrule
\end{tabular}
\end{center}

\paragraph{Cuándo conviene cada uno.} XML sigue siendo preferible cuando
hay que validar con un esquema estricto, transformar a varios formatos
de salida o comunicarse con sistemas legados. JSON es la elección por
defecto para APIs nuevas, almacenamiento documental y cuando el
consumo es código JavaScript del lado del cliente.

% =====================================================
\section{Cómo se conecta con el Práctico 7}
% =====================================================

\begin{itemize}
    \item \textbf{Ejercicio 1}: aplicación directa de JsonPath sobre un
    documento JSON (operadores de navegación y filtros).
    \item \textbf{Ejercicio 2}: generación programática de JSON desde
    Java con Gson, partiendo de datos relacionales.
    \item \textbf{Ejercicio 3}: modelado documental con embedding
    (\texttt{direccion}, \texttt{telefonos}) y referencing
    (\texttt{noticias}, \texttt{comentarios}).
    \item \textbf{Ejercicio 4}: operaciones de inserción con
    \texttt{upsert} para idempotencia.
    \item \textbf{Ejercicio 5}: consultas \texttt{find} con proyección y
    \texttt{aggregate} con \texttt{\$lookup} para join lógico entre
    colecciones.
\end{itemize}

% =====================================================
\section{Conclusión}
% =====================================================

Las ideas centrales a retener:

\begin{itemize}
    \item XML y JSON son los formatos de datos semiestructurados que
    dominan el intercambio entre aplicaciones. JSON está reemplazando a
    XML por su simpleza y por su buena integración con APIs REST y
    bases NoSQL.
    \item XML tiene un ecosistema maduro de esquemas (DTD, XML Schema)
    y consultas (XPath, XQuery, XSLT). JSON ofrece equivalentes más
    livianos: JSON Schema y JsonPath.
    \item Las bases documentales (MongoDB) usan JSON como modelo de
    datos. El modelado deja de pivotear sobre la primera forma normal y
    pasa a depender de decisiones de \emph{embedding} vs
    \emph{referencing} según el uso esperado.
    \item Las consultas en MongoDB se hacen con \texttt{find}
    (filtro + proyección) o con el \emph{aggregation framework}
    (pipeline de etapas), que cubre desde joins lógicos con
    \texttt{\$lookup} hasta agregaciones grupales.
\end{itemize}

\end{document}
